\documentclass[11pt]{article}
\usepackage[utf8]{inputenc}
\usepackage[T1]{fontenc}
\usepackage[margin=1in]{geometry}
\usepackage{booktabs}
\usepackage[hidelinks,colorlinks=false]{hyperref}
\usepackage{amsmath}
\usepackage{listings}
\usepackage{xcolor}
\usepackage{array}
\usepackage{longtable}
\usepackage{caption}
\usepackage{microtype}
\usepackage{lscape}
\usepackage{tabularx}
\usepackage{setspace}
\usepackage{parskip}

\lstset{
  basicstyle=\small\ttfamily,
  breaklines=true,
  frame=single,
  backgroundcolor=\color{gray!10},
  columns=flexible,
  keepspaces=true,
}

\captionsetup{font=small, labelfont=bf}

\title{\textbf{Can LLM Agents Test Social-Science Theory?}\\[0.4em]
A Confound-Aware Framework and a Seven-Generation\\
Negative Result on Bloom's 2-Sigma Problem in LLM Agents}

\author{
  Masumi Kawasaki\\
  Independent researcher\\
  \texttt{geeknees@gmail.com}
}

\date{2026-06-16}

\begin{document}

\maketitle

\noindent\textbf{Status:} Working paper --- not peer-reviewed.\\
\textbf{License:} MIT License (see LICENSE in repository root).\\
\textbf{Preprint category (intended):} cs.AI, cs.MA\\
\textbf{Artifact:} Full code, configs, version-by-version analyses, per-run SQLite databases, and raw run directories are released with this paper (the ablation run's database was lost; its generated run report is released verbatim instead). The released analysis documents carry all reported numbers and run identifiers, and the headline figures of every production run have been re-verified against the released databases.\\[0.5em]
\textbf{AI involvement statement:} This manuscript and the underlying research artifact --- the experiment code, the design iterations across all generations, the analyses, and the paper text itself --- were produced end-to-end by LLM agents (Anthropic Claude; the experiments and drafting principally used \texttt{claude-sonnet-4-6}), operating under the direction of the human author, who framed the research question, approved each generation's design, and reviewed the outputs. The human author takes full responsibility for the content. The same model family also serves as the experimental subjects; we discuss this reflexivity in \S9.

\hrule\vspace{1em}

% -----------------------------------------------------------------------
\begin{abstract}
Large language model (LLM) agents are increasingly proposed as simulated subjects
for social-science research: cheaper, faster, and more controllable than human
participants.  We ask whether such agents can serve as a genuine \emph{test bed} for
an educational theory --- Benjamin Bloom's ``2-sigma problem,'' the claim that
one-to-one mastery tutoring produces learning gains roughly two standard deviations
above conventional group instruction.  Over seven design generations (v3--v9c2) we
operationalize ``learning'' as the formation of condition-specific compact memory
from an education-phase transcript, evaluated on a synthetic rule domain
(\emph{Zarn Tokens}) under a strict memory-only protocol that prevents the model
from answering from prior knowledge.

We report two contributions.  First, a \textbf{negative result for LLM agents}:
across every generation, group instruction met or exceeded one-to-one tutoring, and
in our most rigorous experiment (v9c2; 154 learners, 2.4M tokens) the spread across
pedagogical conditions collapsed to 10 percentage points while \emph{learner-type}
differences spanned 38 points.  Bloom's social-interaction advantage did not robustly
reproduce in this agent environment; what looked like a large discussion effect in an
earlier generation (v9c, 37-point spread) is, we argue, largely attributable to an
uncontrolled prerequisite-knowledge confound (we are candid in \S8--\S9 that this
attribution rests on a single before/after comparison, not an ablation).  Second, and
more durable, a \textbf{confound-aware methodology}: an independent audit of our own
pipeline surfaced five families of confound (F1--F5, one of which decomposes into
sub-confounds, for seven specific threats in total) --- static corrective feedback, an
exact-match scoring artifact, missing context injection into discussion agents,
pseudoreplication, and single-instance role agents plus learning-budget asymmetry ---
that silently inflated effect sizes in agent-based theory testing.  We give a
remediation recipe (A0/A3/A4/A5/A6/A8), including a \emph{population-multiplication}
design for true replication and a \emph{readiness gate} that controls prerequisite
knowledge before the manipulated phase.  We argue these confounds are not specific to
our study but are structural risks for the broader program of using LLM agents to
test social-scientific claims, and that negative results obtained under such controls
are themselves informative.
\end{abstract}

\hrule

% -----------------------------------------------------------------------
\section{Introduction}

A growing literature treats LLMs as \emph{silicon subjects} --- generative agents
that can stand in for human participants in economics games, opinion surveys, social
simulations, and pedagogy.  The appeal is obvious: an agent population can be
instantiated in minutes, perfectly logged, and re-run under counterfactual
manipulations that would be unethical or impossible with people.  If LLM agents
reproduce known human findings, they promise a fast loop for theory exploration; if
they diverge, the divergence itself may be diagnostic.

This promise rests on an assumption that is rarely stress-tested: that an agent
experiment, once built, actually measures the construct it claims to measure.  Our
paper is a sustained attempt to falsify that assumption on a single, well-specified
theory.

We chose Bloom's \textbf{2-sigma problem} \citep{bloom1984} as the target.  Bloom
reported that students tutored one-to-one under mastery-learning conditions performed
about two standard deviations above students in conventional 30-to-1 classrooms --- a
gap large enough that the bottom of the tutored distribution roughly matched the top
of the conventional one.  The finding has shaped decades of educational technology
research and is frequently invoked today to motivate AI tutors.  It is also a
\emph{structural} claim --- about class size, individualization, and active
engagement --- which makes it a natural candidate for agent-based operationalization:
the structural variables (who talks to whom, how large the group, how much each
learner participates) are exactly the variables an agent harness can manipulate
cleanly.

Our central question is therefore two-layered:
\begin{enumerate}
  \item \textbf{Substantive:} When the structural conditions of Bloom's theory are
    operationalized as LLM-agent education sessions, do tutoring/small-group
    conditions produce measurably better downstream problem-solving than large-group
    or lecture conditions?
  \item \textbf{Methodological:} Can we trust the answer?  What does it take to make
    an agent-based theory test internally valid?
\end{enumerate}

We answer the first question in the negative and treat the second as our main
contribution.  Over seven design generations we repeatedly found that group
instruction matched or beat one-to-one tutoring --- the opposite of Bloom's
prediction --- and that the apparent effect sizes were extremely sensitive to
methodological choices that have nothing to do with pedagogy.  An independent audit
of our own code uncovered five confound families that had inflated earlier results.
After remediating them, the most rigorous experiment showed pedagogical condition to
be a minor factor (10-point spread) dwarfed by the \emph{cognitive profile} of the
learner agent (38-point spread).

We make the following contributions:
\begin{itemize}
  \item \textbf{A reusable paradigm} for agent-based theory testing in which learning
    is operationalized as memory formation under a memory-only evaluation protocol on
    a synthetic domain immune to prior knowledge (\S4).
  \item \textbf{A seven-generation case study} documenting how design decisions
    changed conclusions, with all run identifiers, sample sizes, and token costs
    reported for reproducibility (\S5).
  \item \textbf{A confound taxonomy (F1--F5) and remediation recipe (A-codes)} for
    agent experiments, including the \emph{population-multiplication} design that
    converts pseudoreplicated single discussions into genuinely replicated ones, and a
    \emph{readiness gate} that controls prerequisite knowledge (\S6).
  \item \textbf{A confound-controlled negative result} on Bloom's 2-sigma problem in
    this environment, plus the finding that learner-profile heterogeneity dominates
    pedagogical condition (\S7--\S8).
\end{itemize}

We are explicit about scope (\S9): we do not claim LLM agents are human learners,
that Zarn Tokens resemble any real curriculum, or that Bloom's theory is false for
humans.  We claim something narrower and, we think, more useful: that the \emph{naive}
agent-based test of this theory yields a strong but spurious effect, and that
obtaining a trustworthy null requires controls that the agent-simulation literature
does not yet treat as standard.

% -----------------------------------------------------------------------
\section{Related Work}

\paragraph{LLMs as simulated human subjects.}
Recent work uses LLMs to emulate survey respondents, economic-game players, and
populations of social agents \citep{argyle2023, aher2023}, and to build ``generative
agent'' societies whose emergent behavior is studied as data \citep{park2023}.  A
parallel strand cautions that such simulations can exhibit \emph{caricatured} or
distributionally narrow behavior, that prompt and sampling choices drive results, and
that apparent realism can mask the absence of the underlying mechanism.  Our work
contributes a concrete, auditable instance of the latter caution: an agent experiment
that produced a large, theory-consistent-looking effect which dissolved under
replication and confound control.

\paragraph{Bloom's 2-sigma problem and its modern reception.}
Bloom's \citeyearpar{bloom1984} original claim motivated the ``search for group
methods as effective as one-to-one tutoring'' and is routinely cited in
intelligent-tutoring-systems and AI-in-education work.  Subsequent human research has
complicated the 2-sigma figure (effect sizes for tutoring are typically smaller, and
mastery learning, feedback, and time-on-task each contribute).  We do not adjudicate
the human literature; we use the theory's \emph{structural} form as a target that
agent harnesses can manipulate.

\paragraph{Internal validity in computational experiments.}
Pseudoreplication --- treating non-independent measurements as independent replicates
--- is a classic threat in ecology and the behavioral sciences \citep{hurlbert1984}
and, we argue, a pervasive and under-recognized threat in agent simulations, where it
is tempting to generate one interaction transcript and let many ``learner'' reads of
it act as a sample.  Likewise, single-instance ``role'' agents (one teacher, one
judge) confound the manipulated variable with the idiosyncrasies of a particular
generated persona.  Our F1--F5 taxonomy collects these and other threats into a
checklist tailored to LLM-agent pedagogy experiments.

\paragraph{Evaluation artifacts in LLM scoring.}
Exact-match grading of free-text model output is known to undercount semantically
correct answers \citep{rajpurkar2016}.  We document a concrete case (our ``L6''
inductive task scoring 0\% for three generations) in which an exact-match scorer, not
the agents, produced a null, and we discuss when a semantic judge is required.

% -----------------------------------------------------------------------
\section{Bloom's Theory, Operationalized}

Bloom's argument has three structural levers that an agent experiment can manipulate
independently:
\begin{enumerate}
  \item \textbf{Group size} --- one tutor's attention divided among 1, 2, \ldots, N
    learners.
  \item \textbf{Individualization} --- whether instruction adapts to the specific
    learner's errors (diagnose $\to$ correct $\to$ retest) or is delivered uniformly.
  \item \textbf{Active engagement / ownership} --- whether a learner is ``called on,''
    contributes, and receives feedback, versus passively receiving a lecture.
\end{enumerate}

Bloom predicts that small group size, high individualization, and high engagement
jointly produce the 2-sigma gain.  Our generations probe these levers in different
combinations: early generations (v4--v6) contrast a 1:N classroom against 1-on-1
tutoring (levers 1+2); the mechanism tests (v7) isolate engagement vs.\
individualization; the later generations (v8--v9c2) hold a common lecture fixed and
vary group size and ownership in the \emph{post-lecture} phase (levers 1+3), which
most directly matches a ``flipped classroom'' reading of mastery learning.

A crucial operational decision is what counts as ``learning.''  We do \textbf{not}
fine-tune weights.  We operationalize learning as the creation, during an education
phase, of a compact \textbf{memory} object that the agent must rely on at evaluation
time --- the rules are withheld from the evaluation prompt (\S4.2).  Condition
differences thus manifest as differences in the \emph{content and structure of memory}
that each pedagogical format induces, and in downstream task accuracy that depends on
that memory.  This is a deliberately modest construal of learning, but it is the one
an agent harness can measure without confounding ``learning'' with ``rule reading.''

% -----------------------------------------------------------------------
\section{The Zarn Token Paradigm}

\subsection{A synthetic domain immune to prior knowledge}

To prevent the model from solving tasks from pre-training, we invented a closed
scoring system, \textbf{Zarn Tokens}, with no presence in any training corpus.  A
sequence of colored tokens is scored by four interacting rules plus a stacking rule:

\begin{table}[h]
\centering
\caption{Zarn Token rules.}
\label{tab:zarn-rules}
\begin{tabular}{@{}llp{7cm}@{}}
\toprule
Token & Base value & Rule \\
\midrule
Red    & 3 (but see below) & \textbf{Modifier only.} Red's own contribution is always 0; it doubles the base value of the token immediately to its right. \\
Blue   & 5 & \textbf{Conditional.} Active (5) only if at least one Green appears \emph{anywhere} to its left; otherwise 0. \\
Green  & 2 & \textbf{Positional.} Active (2) in any position \emph{except} last; inactive (0) if it is the final token. \\
Yellow & 7 & \textbf{Always active} (7), regardless of position or neighbors. \\
\bottomrule
\end{tabular}
\end{table}

\textbf{Stacking rule:} doubling an \emph{inactive} token still yields 0 (e.g.,
\texttt{[Red, Blue]} with no Green scores 0, not 10).  The intended evaluation
procedure is \emph{activation-check first, then modifier application} --- a
procedural-ordering requirement that becomes a probe for one of our learner types
(\S4.3).

The domain is small enough to teach in one lecture yet rich enough to support a
difficulty ladder: conditional activation (Blue), positional inactivation
(Green-at-end), modifier composition (Red doubling), and the interaction of all three
(doubled-inactive).  Worked example: \texttt{[Green, Red, Yellow, Blue]} $\to$ Green
active (2); Red modifies Yellow; Yellow doubled (14); Blue active because Green is to
its left (5) $\to$ \textbf{21}.

\subsection{Memory-only evaluation protocol}

Each evaluation task carries two fields: a \texttt{learner\_prompt} with the rules
\textbf{removed}, and \texttt{hidden\_rules} used \textbf{only} by the auto-scorer.
The learner must answer from memory formed during the education phase, with an
explicit instruction \emph{``Use only your memory of the rules.  Do not assume rules
you were not taught.''}  This split is the methodological heart of the paradigm.  In the earliest pilot
(v1), evaluation tasks never contained the rules text, but memory formation was
unconstrained (no word budget, no rule-exposure limit): the teacher's lecture
recited the domain rules near-verbatim, the transcript carried them into memory
largely intact, and the solver then answered from a memory that was effectively
a copy of the rulebook --- every condition scored 100\% (a complete ceiling) by
a different route than we originally assumed.  In v2 the rules were additionally
embedded directly in the evaluation-task prompt text, which reproduced the same
ceiling through a second, more direct channel.  Both routes mask pedagogical
effects.  Removing the rules from the evaluation prompt (v3 onward) and
separating \texttt{learner\_prompt} from \texttt{hidden\_rules} dropped the
no-education baseline to (near) 0\% and opened measurable variance.

\subsection{Learner types as post-generation memory constraints}

The memory object is a structured artifact the model produces from the education
transcript.  It has seven fields:

\begin{lstlisting}[language={}]
{
  "rules": [], "examples": [], "edge_cases": [], "strategy": [],
  "corrected_misconceptions": [], "remaining_misconceptions": [], "uncertain_rules": []
}
\end{lstlisting}

At evaluation time this object --- not the transcript --- is what the learner reasons
from.  ``Learning differences'' between conditions are therefore differences in the
content of these fields.

A recurring early failure (v5) was that \emph{prompt-level} learner personas
(``low-ability,'' ``forgets edge cases'') barely changed behavior, because the
underlying model capability is fixed.  From v6 onward we instead impose
\textbf{hard constraints on the memory object after the model generates it}, yielding
theory-driven learner types:

\begin{table}[h]
\centering
\caption{Learner-type memory constraints applied after memory generation.}
\label{tab:learner-types}
\begin{tabular}{@{}lrrrr@{}}
\toprule
Type & Memory budget & Edge-case retention & Rule-order retention & Asks questions \\
\midrule
\texttt{rule\_extractor}    & 150 words & 90\% & 90\% & 30\% \\
\texttt{edge\_case\_dropper} & 120 words & 10\% & 70\% & 20\% \\
\texttt{order\_confused}     & 100 words & 60\% & 20\% & 40\% \\
\texttt{passive\_listener}   &  80 words & 20\% & 40\% &  0\% \\
\bottomrule
\end{tabular}
\end{table}

(Three further types --- \texttt{help\_seeker}, \texttt{answer\_first},
\texttt{example\_memorizer} --- were defined in v6 and used in early generations.)
Constraints are applied mechanically: each edge-case memory item is dropped
independently with probability $1 - \text{retention}$; with probability $1 -
\text{rule\_order\_retention}$ the ordered rule list is shuffled (degrading
procedural knowledge); memory is then trimmed to the word budget by dropping
lowest-priority fields first; and a \texttt{passive\_listener} with question
probability 0 never speaks in discussion (its turn is skipped with no model call).
These constraints make learner heterogeneity \emph{real} at the representation level
rather than role-played.

\subsection{Phase structure}

The mature protocol (v8 onward) runs four phases:
\begin{enumerate}
  \item \textbf{Lecture (Phase 1).}  A common instructional artifact.  Early
    generations let the model generate the lecture per run; from v9c we fix a single
    rich reference lecture (6,493 characters, five worked examples, an edge-case
    checklist, and a common-mistakes table) so that lecture quality is held constant
    across conditions.
  \item \textbf{Readiness check (Phase 2).}  A short diagnostic (3 items in v8; 4
    types --- recall, edge\_case, rule\_interaction, procedure\_order --- in v9c+)
    that auto-scores each learner and appends a corrective note to memory on wrong
    answers.  From v9c we add a \textbf{readiness gate}: if the cohort does not reach
    an 80\% pass rate, a \texttt{readiness\_failed} flag is set, signaling that
    prerequisite knowledge was \emph{not} established and that any downstream
    condition contrast is confounded.
  \item \textbf{Condition-specific interaction (Phase 3).}  The manipulated phase:
    lecture-only vs.\ discussion of various group sizes, or tutoring, depending on
    the generation.
  \item \textbf{Evaluation (Phase 4).}  Ten memory-only tasks spanning a Bloom-style
    ladder.
\end{enumerate}

\subsection{Tasks, scoring, and metrics}

The evaluation ladder has ten tasks: \texttt{recall} (L1), \texttt{edge\_case}
$\times$2 (L2), \texttt{rule\_interaction} (L3), \texttt{debugging} $\times$3
(L4/L4/L5), \texttt{short\_rule\_induction} (L6), \texttt{peer\_error\_detection},
and \texttt{explanation\_choice}.  Numeric and multiple-choice answers are graded by
\textbf{exact match} against a hidden answer; this is reliable for the nine
closed-form tasks but, as \S6 details, fails badly for the free-text inductive task
(L6), which we ultimately exclude from the main score and re-route to a semantic
judge.

Beyond accuracy we log: an \textbf{ownership} score (contributions, share called-on,
share receiving feedback), a \textbf{misconception-correction} rate (did the
diagnosed misconception disappear from later memory?), a 10-item
\textbf{memory-coverage} vector and its pre/post-discussion \textbf{delta},
\textbf{confidence calibration} (high-confidence-wrong rate), and full \textbf{token
accounting} per phase.  Token accounting matters because tutoring and large-discussion
conditions are far more expensive than lecture-only, and a fair comparison must report
cost as well as accuracy.

% -----------------------------------------------------------------------
\section{The Experimental Arc (v3--v9c)}

We summarize the production generations v4--v9c below; v3 was the pilot that
established the memory-only protocol (\S4.2), and the confound-controlled v9c2 is
treated separately in \S7.  (The ``seven generations'' of the title counts design
generations --- v3, v4, v5, v6, v7, v8, and the v9 series --- with v7a/v7b as two
experiments within one generation and v9b/v9c/v9c2 as iterations of one.)  All
production runs use
\texttt{claude-sonnet-4-6} for every agent role (one v9b smoke run used
\texttt{claude-haiku-4-5}).  The no-education baseline scored 0\% in every
production run (v4 onward), confirming that the memory-only design removes the
prior-knowledge ceiling; we omit it from the tables below except where noted.
(In one earlier v3 pilot smoke test, before the production protocol was frozen,
a no-education learner scored 1/8 by correctly judging a since-retired L5
counterexample-format task; post hoc inspection showed that task's true/false
judgment was answerable by pure logical inference on the claim's structure,
without any domain-rule knowledge, so it was replaced with the debugging-format
L5 used in all production runs (\S4.5).  This pilot-era artifact does not
affect any production run reported below.)

\subsection{v4--v6: classroom vs.\ tutoring, and the homogenization effect}

\begin{table}[h]
\centering
\caption{v4--v6 summary.}
\label{tab:v4-v6}
\begin{tabular}{@{}llp{6cm}p{4.5cm}@{}}
\toprule
Gen & run\_id (prefix) & Conditions and main accuracy & Headline \\
\midrule
v4 & \texttt{19b39e20} & no-ed 0\%, \textbf{classroom 47\%}, tutoring 28\% ($n$=3/4/4) & Bloom reversed: classroom $>$ tutoring \\
v5 & \texttt{02070b48} & homogeneous-classroom 44\%, heterogeneous-classroom 41\%, \textbf{tutoring 38\%} ($n$$\approx$4 each) & Tutoring lowest; highest variance \\
v6 & \texttt{cd5d2f0e} & homogeneous-classroom \textbf{94\%}, heterogeneous-classroom 59\%, tutoring 59\% ($n$=4 each) & Format effect $\approx$ 0 (hetero = 1on1); homogenization effect appears \\
\bottomrule
\end{tabular}
\end{table}

Three findings recur and survive into the later generations.
\begin{enumerate}
  \item[(i)] \textbf{Bloom's tutoring advantage never appears}: classroom matches or
    beats one-to-one in every generation.
  \item[(ii)] A \textbf{homogenization effect}: in v6, relative to a heterogeneous
    classroom, one-to-one tutoring raised the weakest learner type
    (\texttt{passive\_listener} +25 points) and lowered the strongest
    (\texttt{order\_confused} $-$25 points), compressing variance --- tutoring made
    the group more uniform without raising the mean.  (The variance-compression
    direction is consistent with v5, where tutoring also showed the highest
    per-learner spread; the $\pm$25-point figure itself is a v6 $n$=4-per-type result
    and should be read as a hypothesis, not an estimate.)
  \item[(iii)] \textbf{Misconception correction is decoupled from score}: v6 tutoring
    corrected diagnosed misconceptions 75\% of the time yet did not outscore the
    classroom, because correcting a \emph{propositional} misconception (``Blue is
    always active'') does nothing for a \emph{procedural} deficit (shuffled rule
    order).  This decoupling reappeared in v7b (generic tutoring: 100\% correction,
    47\% accuracy) and v8.
\end{enumerate}

The v6 ``homogeneous classroom 94\%'' is the one apparent win for group instruction,
but it is a confounded one: that condition's teacher was told the whole class shared
a single misconception and to target it --- an efficiency artifact of asymmetric
design, not a clean format comparison.

\subsection{v7: mechanism tests}

Two targeted experiments tried to explain the v6 cross-over.

\begin{itemize}
  \item \textbf{v7a (passive\_listener rescue; \texttt{dd2fd888}):}  Does tutoring
    help the passive learner via \emph{forced interaction} or via
    \emph{personalization}?  Conditions: public-Q\&A lecture (88\%), forced-check-in
    lecture (53\%), one-to-one (50\%).  Forced-check-in $\approx$ one-to-one (3
    points) suggested the benefit is contact, not personalization --- but a
    lecture-length asymmetry (a stray ``under 200 words'' constraint shortened the
    check-in lecture to 1,022 vs.\ 3,599 characters) confounds the comparison, and
    the richest \emph{passive} lecture scored highest.
  \item \textbf{v7b (order\_confused intervention; \texttt{32ffaa49}):}  Does
    procedure-scaffolded tutoring beat generic tutoring for the procedurally confused
    learner?  Conditions: public-Q\&A (69\%), generic tutoring (47\%),
    procedure-scaffolded tutoring (41\%).  Scaffolding did not help and slightly hurt;
    generic tutoring reached 100\% misconception correction yet 47\% accuracy ---
    again, correction decoupled from outcome.  A bimodal \texttt{order\_confused}
    distribution (one learner at 0/8) made $n$=4 means unstable.
\end{itemize}

The v7 mechanism tests thus reinforced the negative result but also exposed how
fragile $n$=4 condition means are and how easily an incidental prompt difference
(lecture length) masquerades as a pedagogical effect --- foreshadowing the formal
confound audit.

\subsection{v8--v9c: flipped learning, class size, and the readiness problem}

From v8 we adopt the flipped design: a \emph{common} lecture for all conditions, then
a manipulated post-lecture phase.

\begin{table}[h]
\centering
\caption{v8--v9c summary.}
\label{tab:v8-v9c}
\begin{tabular}{@{}llrp{5.5cm}p{4cm}@{}}
\toprule
Gen & run\_id (prefix) & $n$ & Main accuracy by condition & Headline \\
\midrule
v8  & \texttt{ef107ef9} & 16 & whole-class-discussion \textbf{83\%}, one-to-one 78\%, lecture-only 70\%, small-group 68\% & Discussion $\geq$ lecture; but lecture length still varied by condition \\
v9b & \texttt{772d2ce7} & 34 & lecture-only \textbf{60\%} = pair \textbf{60\%}, small 48\%, large 47\%, medium 45\% & Class size \textbf{not} monotone; ownership--score correlation $\approx$ 0 \\
v9c & \texttt{b412cfdb} & 34 & pair \textbf{65\%}, small 57\%, large 52\%, medium 40\%, lecture-only \textbf{28\%} & 37-point spread; ``discussion adds value'' --- \emph{but readiness failed (69\%)} \\
\bottomrule
\end{tabular}
\end{table}

v9b introduced the class-size design (instantiating 2/4/8/16-learner classes; not
to be confused with the \emph{population-multiplication} replication design of
\S6, which arrives only in v9c2) and an explicit ownership metric, and found neither a
monotone class-size effect nor an ownership--outcome correlation: lecture-only tied
the smallest discussion for the top score, and memory ``delta'' from discussion was
$\sim$0 (agents acquired essentially no new rule knowledge from discussing).  v9c
fixed the lecture and added a four-type readiness check --- and produced the most
Bloom-friendly numbers in the entire program: a 37-point spread with every discussion
condition beating lecture-only (28\%).  It would have been tempting to stop here and
report a discussion advantage.

We did not, because the readiness gate had \textbf{failed}: the cohort reached only
69\% on the readiness check (recall 68\%, rule\_interaction 50\%), below the 80\%
target.  The condition contrast was therefore measured over learners who had \emph{not}
reliably acquired the prerequisites --- exactly the regime in which downstream
differences reflect noise in prerequisite acquisition rather than the manipulated
phase.  The lecture-only condition, moreover, had \emph{zero} learning sessions in
Phase 3 while discussion conditions had one, a raw learning-budget asymmetry.
Notably, the \emph{learner-type} spread in v9c was only 8 points
(\texttt{rule\_extractor} 52\% vs.\ \texttt{passive\_listener} 44\%) --- far below
the 37-point \emph{condition} spread, the inverse of the controlled v9c2 ordering
(\S7) --- which is itself a signature that the v9c condition differences were not
stable learner-driven effects.  The striking v9c effect was a prime suspect for
confounding, which motivated a formal audit.

% -----------------------------------------------------------------------
\section{Confound Taxonomy and Remediation}

We commissioned an independent audit of our own pipeline (a fresh analysis pass over
code and per-run databases, using two scripts --- \texttt{check\_replication.py} and
\texttt{check\_confounds.py} --- that we release).  It produced a cross-version audit
table and five confound families (F1--F5; F5 decomposes into F5a/b/c).  We present
them as a reusable checklist for agent-based experiments, with the remediation codes
(A-codes) we implemented.  The A-codes are the project's internal change-IDs from the
v9c2 implementation plan, not a re-derived sequence; they are therefore
non-contiguous (A0, A3, A4, A5, A6, A8) and do \textbf{not} align numerically with
the F-codes.  Appendix B gives the explicit confound$\to$remediation map; we retain
the original IDs so the paper matches the released code rather than renaming for
cosmetic alignment.

\paragraph{F1 --- Static corrective feedback.}
Wrong readiness-check answers append a \emph{fixed, pre-written} corrective note to
memory, identical regardless of which misconception caused the error.  This does not
bias the pass/fail score (scoring precedes the note) but undermines any claim that
``adaptive feedback'' drove learning.  \emph{Status:} acknowledged design limitation;
not fully remediated (LLM-generated per-error notes are the proper fix and remain
future work).

\paragraph{F2 --- Exact-match scoring artifact $\to$ A8.}
The free-text inductive task (L6, \emph{short\_rule\_induction}) scored 0\% in v8,
v9b, and v9c.  Re-scoring with curated aliases recovered almost nothing and did not
generalize, but inspection confirmed the agents were producing \emph{semantically
correct} rule statements that the exact-match scorer rejected.  This is an evaluation
artifact, not an agent failure.  \emph{Remediation A8:} exclude L6 from the main
score and route it to an LLM semantic judge (\texttt{SEMANTIC\_TASK\_TYPES}),
implemented for subsequent runs.  (The nine closed-form tasks are unaffected and
remain exact-match.)

\paragraph{F3 --- Discussion context not injected $\to$ A0.}
In the discussion phase, \emph{participant} agents received neither the lecture text
nor their own prior memory in their prompt --- only the moderator did.  Transcripts
confirmed called-on learners literally opening with ``I don't actually know the
rules.''  Between-condition differences therefore partly measured \emph{how an agent
improvises from an empty context}, not collaborative knowledge construction.
\emph{Remediation A0:} inject each learner's memory into the contribution and reply
prompts so that participants reason from what they know.

\paragraph{F4 --- Pseudoreplication $\to$ A3.}
This is the most consequential confound.  In v4--v9c, each discussion/classroom
condition generated \textbf{one} transcript (\texttt{unique\_discussions = 1}); the
multiple ``learners'' in a condition were all readers of that single conversation.
The per-condition standard deviation therefore reflected learner-response variance to
one $N$=1 event, \emph{not} uncertainty about the pedagogical design --- yet it is
exactly the quantity one would (wrongly) use to claim ``condition A $>$ condition B.''
\emph{Remediation A3:} a \textbf{population-multiplication} design that generates
\texttt{n\_disc} (=5) \emph{independent} discussions per condition, each with freshly
sampled learners, and aggregates at the discussion level.  This converts
pseudoreplicates into genuine replicates and is, we believe, the single most
important control for agent-based social experiments.

\paragraph{F5a --- Learner-profile homogeneity $\to$ A4 (already satisfied).}
Whether the learner population is diverse.  v7a/v7b deliberately used a single type
(a scoped choice, not a confound); v8/v9b/v9c already used four balanced types.
\emph{Status:} satisfied from v8 on; an erroneous audit note claiming ``all learners
were one type'' was traced to a smoke-config database and corrected.

\paragraph{F5b --- Single-instance role agents $\to$ A5.}
One classroom teacher, one tutor, one evaluator instance served every condition in
all of v4--v9c.  The manipulated variable is thus confounded with the idiosyncrasies
of a single generated persona/judge.  \emph{Remediation A5:} instantiate a fresh
teacher agent for each of the \texttt{n\_disc} independent discussions (five distinct
teachers).  We deliberately did \textbf{not} triplicate the evaluator: because the
nine retained tasks are exact-match graded, the scoring path never reaches an LLM
judge, so evaluator multiplication would be a no-op for them (it becomes relevant
only once the L6 semantic judge is in use).

\paragraph{F5c --- Learning-budget asymmetry $\to$ A6.}
Lecture-only learners had zero Phase-3 sessions while discussion learners had one, so
``lecture-only vs.\ discussion'' confounded \emph{format} with \emph{amount of
post-lecture learning opportunity}.  \emph{Remediation A6:} a \textbf{self-reflection
phase} giving lecture-only learners a structured individual-reflection session --- an
\emph{additional learning opportunity} matched in kind (though, as we report
candidly, not in token volume; see \S7 and \S9).

The cross-version audit table (released as an artifact) records, for every generation
v4--v9c, the run-id$\leftrightarrow$database mapping, the \texttt{sessions} vs.\
\texttt{unique\_discussions} count per condition (F4), and the single-instance/diversity
status (F5).  Its summary matrix shows F4 present in every generation except v8
(partial) and F5b present in all --- i.e., the two structural confounds were
\emph{baked in from the first experiment} and silently shaped six generations of
results before the audit caught them.  This is the cautionary core of the paper: the
confounds were invisible at the level of ``the code runs and produces plausible
numbers,'' and only a replication-and-independence audit surfaced them.

% -----------------------------------------------------------------------
\section{v9c2: The Confound-Controlled Experiment}

v9c2 applies A0/A3/A5/A6/A8 together (A4 --- learner-profile diversity --- was
verified as already satisfied from v8 on) and re-asks the v9c questions under control.
(run\_id \texttt{dc1bdd27}; \texttt{claude-sonnet-4-6}; 154 learners; 2,376,593
tokens.)

\textbf{Design.}  Five conditions: \texttt{lecture\_plus\_self\_reflection} ($n$=4;
the renamed, A6-corrected lecture-only) and four discussion sizes --- pair (2), small
(4), medium (8), large (16) --- each realized via population-multiplication with
\texttt{n\_disc = 5} independent discussions and a fresh teacher per discussion.
Four balanced learner types.  Common fixed lecture; four-type readiness check with
the 80\% gate.

\textbf{The readiness gate passed.}  The cohort reached \textbf{90\%} readiness
(edge\_case 100\%, recall 91\%, procedure\_order 85\%, rule\_interaction 85\%) --- the
first generation to clear the gate, so the condition contrast is, for the first time,
measured over learners who genuinely held the prerequisites.

\textbf{Main result --- the effect collapses.}

\begin{table}[h]
\centering
\caption{v9c2 accuracy and token efficiency by condition.}
\label{tab:v9c2-main}
\begin{tabular}{@{}lrrl@{}}
\toprule
Condition & Accuracy & Edu tokens & Correct per 1k edu tokens \\
\midrule
\texttt{lecture\_plus\_self\_reflection} & 83\% &  2,716 & 11.05 \\
\texttt{pair\_discussion\_size\_2}        & 80\% & 45,775 &  1.57 \\
\texttt{small\_class\_discussion\_size\_4} & 76\% & 42,590 &  3.19 \\
\texttt{medium\_class\_discussion\_size\_8} & 86\% & 41,580 &  7.46 \\
\texttt{large\_class\_discussion\_size\_16} & 81\% & 35,035 & 16.67 \\
\bottomrule
\end{tabular}
\end{table}

The spread across pedagogical conditions is \textbf{10 percentage points} (medium
86\% vs.\ small 76\%), down from v9c's 37 --- and the ordering is non-monotone in
both class size and ownership.  The interpretation flags computed by the harness are:
\texttt{class\_size\_effect\_supported = false}, \texttt{ownership\_effect\_supported
= false}, \texttt{discussion\_added\_value = false}.  We stress that these flags are
\emph{threshold} tests (an effect must exceed a fixed margin over the lecture
baseline), not significance tests; the honest reading is ``no \emph{large} effect,''
not ``exactly zero.''  Indeed the numerically top condition is a discussion condition
(medium, 86\%), one point above lecture-plus-self-reflection (83\%) --- so we do not
claim discussion \emph{hurts}, only that no condition clears a meaningful margin and
the v9c ``discussion advantage'' (every discussion condition far above a 28\% lecture
baseline) did \textbf{not} survive readiness control and replication.

\textbf{Learner type dominates.}  The spread across learner \emph{types} is
\textbf{38 points} (\texttt{edge\_case\_dropper} 96\% vs.\ \texttt{passive\_listener}
58\%) --- nearly four times the condition spread.  On the same $n$ that makes us
cautious about the condition contrast, this larger and directionally stable gap (the
\texttt{passive\_listener}, the most memory-constrained type, is lowest in v6, v9b,
v9c, and v9c2) is the more defensible comparative claim: the single largest
determinant of downstream accuracy is the cognitive profile (the memory constraint),
not the pedagogical format.  The v9c learner-type spread had been only 8 points; with
the larger, readiness-controlled sample the profile differences resolve far more
clearly.  We still report this as a spread, not a tested effect (\S9).

\textbf{Token efficiency reframes ``what won.''}  Because A6 matched lecture-only's
learning opportunity in \emph{kind} but not in \emph{volume},
\texttt{lecture\_plus\_self\_reflection} spent $\sim$2.7k education tokens against
35--46k for discussion conditions.  Per-1k-token learning efficiency therefore ranks
lecture-only (11.05) and large-class (16.67) far above pair/small discussion
(1.57/3.19).  Two caveats apply.  Large-class scores highest on this metric only
because one teacher's tokens are amortized over 16 learners --- this is an
instructional-cost-per-head statement, not a learning-quality one, and should not be
read as endorsing large classes.  And the same metric exposes the principal remaining
confound in v9c2 (\S9): A6 achieved ``equal learning opportunity'' structurally (a
reflection session) but not in compute ($\sim$1/15th the tokens), so lecture-only's
83\% is either impressively cheap or unfairly under-resourced depending on which
budget one holds fixed.  The robust statement is the weaker one: across a 1--2
order-of-magnitude range of token spend, downstream accuracy moved by at most 10
points.

% -----------------------------------------------------------------------
\section{Cross-Cutting Findings}

\begin{enumerate}
  \item \textbf{The education effect is real; the social-interaction effect is not.}
    No-education baselines sat at 0\% throughout, so the pipeline does measure
    learning.  But across seven generations the \emph{Bloom-specific} prediction ---
    tutoring/small-group $>$ large-group/lecture --- never robustly reproduced; the
    strongest pro-Bloom signal (v9c) was an artifact of an uncontrolled readiness
    confound.

  \item \textbf{The headline numbers are best explained by confounds, not pedagogy
    --- though the comparison is not a clean ablation.}  The v9c$\to$v9c2 contrast
    moved the condition spread from 37 to 10 points and flipped three interpretation
    flags.  We are deliberately careful about the inference: v9c2 changed six
    remediations \emph{at once} and raised $n$ from 34 to 154, so we cannot attribute
    the collapse to any single control.  The two mechanisms are also partly entangled
    by design --- A3 (replication) inherently increases $n$, and a 37-point spread
    built on \texttt{unique\_discussions = 1} per condition (F4) was, by definition,
    learner-response variance to a single discussion rather than a stable effect.  So
    ``we removed confounds and the effect collapsed'' and ``we added replication and
    $N$=1 noise averaged out'' are, here, two descriptions of the \emph{same}
    remediation rather than rival explanations.  What we can claim firmly is the
    conditional: a spread that does not survive replication and prerequisite control
    was not safe to report as a pedagogical effect.  What we cannot claim is a
    variable-isolated causal decomposition; a single-control ablation is the obvious
    next experiment (\S9).  We did run one such ablation ---
    \texttt{bloom\_v9c2\_ablation\_ndisc1} (run \texttt{a551c74d}; \texttt{n\_disc=1},
    all other v9c2 remediations ON) --- and it produced a 25-point condition spread
    (pair 72\% vs.\ lecture 47\%), intermediate between v9c (37pp) and v9c2 (10pp),
    consistent with replication being a key factor.  However, the ablation's readiness
    gate failed at 53\%, so prerequisite knowledge was also not controlled in this
    arm; the two factors changed simultaneously and the result cannot cleanly isolate
    A3 alone.

  \item \textbf{Learner heterogeneity is the dominant variable.}  When prerequisites
    are controlled, \emph{who the learner is} (its memory constraints) explains far
    more variance (38 points) than \emph{how it was taught} (10 points).  This is
    itself a substantive result: in this environment, instructional format is a
    second-order factor behind learner representation.

  \item \textbf{Two robust micro-mechanisms.}  (i) \emph{Homogenization}: one-to-one
    tutoring compresses the outcome distribution (raising the weakest type, lowering
    the strongest) without raising the mean.  (ii) \emph{Correction-outcome
    decoupling}: high misconception-correction rates do not translate into higher
    accuracy, because correcting a stated misconception leaves procedural deficits
    (and unrelated task demands) untouched.  Both held across multiple generations and
    are candidate hypotheses for human study, stated here as agent-environment
    findings only.

  \item \textbf{Cost asymmetry is large and usually ignored.}  In v9c2, discussion
    conditions spent 13--17$\times$ the education-phase tokens of
    lecture/self-reflection (35--46k vs.\ 2.7k) for equal-or-worse accuracy, and in
    earlier generations one-to-one tutoring cost ${\sim}4.3\times$ the classroom
    lecture (v5: 35,095 vs.\ 8,180 tokens).  Any agent-based pedagogy claim should report
    cost-normalized outcomes; otherwise the ``expensive condition'' can look better
    simply because it did more.
\end{enumerate}

% -----------------------------------------------------------------------
\section{Limitations and Scope}

We constrain our claims tightly.

\begin{itemize}
  \item \textbf{Agents are not students.}  We measure memory formation and rule
    application by a single LLM under imposed memory constraints, not human cognition.
    Nothing here bears directly on whether Bloom's 2-sigma effect holds for people.
  \item \textbf{One domain, one model family.}  All results are on Zarn Tokens with
    \texttt{claude-sonnet-4-6}.  The synthetic domain was chosen for prior-knowledge
    immunity, not ecological validity; generalization across domains and model families
    is untested.
  \item \textbf{Residual token asymmetry (the live confound).}  A6 equalized the
    \emph{kind} of learning opportunity for lecture-only but not its token volume
    ($\sim$2.7k vs.\ $\sim$35--46k).  A cleaner future design would equalize compute
    budget across conditions or report a budget-matched arm.
  \item \textbf{No clean single-variable ablation.}  The pivotal v9c$\to$v9c2
    comparison bundles six remediations and a 4.5$\times$ sample-size increase.  We
    can show the inflated effect did not survive control, but not which control was
    responsible.  We ran one ablation (\texttt{bloom\_v9c2\_ablation\_ndisc1}, run
    \texttt{a551c74d}): \texttt{n\_disc=1} with all other v9c2 remediations ON.  It
    produced a 25-point condition spread --- intermediate between v9c (37pp, readiness
    69\%) and v9c2 (10pp, readiness 90\%) --- directionally consistent with
    replication (A3) being a key driver.  But the ablation's readiness gate failed at
    53\%, meaning prerequisite knowledge was \emph{also} not controlled in this arm.
    Both factors shifted at once, so the ablation confirms that the bundle matters but
    cannot isolate A3 alone.  A cleaner next design would hold readiness fixed (e.g.,
    by running multiple passes until the gate clears) and vary \texttt{n\_disc}
    independently.
  \item \textbf{Statistical power.}  Even with \texttt{n\_disc = 5}, five independent
    discussions per condition do not support strong inferential tests; the 10-point
    condition spread is best read as ``no large effect,'' not ``exactly zero.''  The
    same caution applies in principle to the 38-point learner-type spread --- we lean
    on it more heavily only because it is larger and directionally consistent across
    four generations, not because it is statistically established.  We report
    descriptive statistics and interpretation flags, not $p$-values, and recommend
    \texttt{n\_disc = 10--15} for a confirmatory run.
  \item \textbf{F1 unresolved.}  Corrective feedback remains static; adaptive
    per-error feedback is unimplemented.
  \item \textbf{Stochasticity is not seed-controlled.}  The learner-type constraints
    (probabilistic edge-case dropping, rule-order shuffling) and LLM sampling run
    without fixed RNG seeds, so runs are not bit-reproducible.  Analyses regenerate
    deterministically from each run's frozen outputs, but an exact re-run will
    differ; seeded constraint application is planned for the next generation.
  \item \textbf{Reflexivity of AI authorship.}  As disclosed in the front matter,
    this study was designed, executed, and written by LLM agents of the same model
    family as its experimental subjects, under human direction and review.  We flag
    this as a potential source of shared blind spots --- the auditor and the audited
    share a substrate --- and note that the confound audit (\S6) was itself
    agent-executed; independent human replication of the audit is the obvious
    external check.
  \item \textbf{Construct of ``learning.''}  Memory-only evaluation is a deliberate,
    narrow operationalization; richer constructs (transfer, retention over sessions)
    are out of scope.
\end{itemize}

We view the open items as a strength of the artifact: each is logged, located in
code, and prioritized for the next generation, which is the posture we advocate for
agent-based theory testing generally.

% -----------------------------------------------------------------------
\section{Conclusion}

We set out to test whether LLM agents can serve as a test bed for Bloom's 2-sigma
problem, and answered on two levels.  Substantively, in this environment the theory's
social-interaction advantage does not reproduce: under confound control, pedagogical
format is a minor factor (10-point spread) dominated by learner cognitive profile
(38-point spread), and the social conditions are far more expensive for no accuracy
gain.  Methodologically --- and this is the contribution we expect to outlast the
specific result --- we showed that the \emph{naive} version of the same experiment
produced a large, theory-flattering effect (a 37-point spread, ``discussion adds
value'') that \textbf{did not survive} replication and prerequisite control.  We are
careful not to over-attribute: because the controls were applied together, we
demonstrate that the effect was unsafe to report, not which confound created it.  We
distilled the threats into a five-family confound taxonomy (F1--F5) and a remediation
recipe (A-codes), centered on a population-multiplication design for genuine
replication and a readiness gate for prerequisite control.

The broader lesson for the ``LLM-as-simulated-subject'' program is sobering and, we
hope, constructive: an agent experiment can run cleanly, produce plausible numbers,
and still be measuring the wrong thing.  Negative results obtained \emph{after}
explicit internal-validity controls are more informative than positive results
obtained before them.  We release all code, configurations, version-by-version
analyses, and per-run databases (the ablation run's database was lost; its
generated run report is released verbatim) so that the confound checklist ---
and the cautionary tale --- can be reused.

% -----------------------------------------------------------------------
\appendix

\section{Run Index}
\label{appendix:runs}

\begin{table}[h]
\centering
\caption{All production runs, v4 through ablation.}
\label{tab:run-index}
\footnotesize
\begin{tabular}{@{}llp{3.8cm}lrp{4cm}@{}}
\toprule
Gen & Experiment name & run\_id & Model & $n$ & Notes \\
\midrule
v4       & bloom\_1n\_vs\_1on1\_v4             & \texttt{19b39e20-...} & sonnet-4-6 &  11 & shared experiment.db \\
v5       & bloom\_heterogeneity\_v5            & \texttt{02070b48-...} & sonnet-4-6 &  15 & shared experiment.db \\
v6       & bloom\_learner\_types\_v6           & \texttt{cd5d2f0e-...} & sonnet-4-6 &  15 & shared experiment.db \\
v7a      & (passive\_listener rescue)          & \texttt{dd2fd888-...} & sonnet-4-6 &  12 & \\
v7b      & (order\_confused intervention)      & \texttt{32ffaa49-...} & sonnet-4-6 &  12 & \\
v8       & (flipped learning)                  & \texttt{ef107ef9-...} & sonnet-4-6 &  16 & \\
v9b      & bloom\_v9b (smoke)                  & \texttt{9b599c12-...} & haiku-4-5  &  14 & smoke \\
v9b      & bloom\_v9b (prod)                   & \texttt{772d2ce7-...} & sonnet-4-6 &  34 & 619,507 tokens \\
v9c      & bloom\_v9c\_classroom\_size         & \texttt{b412cfdb-...} & sonnet-4-6 &  34 & readiness\_failed (69\%); 631,955 tokens \\
v9c2     & bloom\_v9c2\_classroom\_size        & \texttt{dc1bdd27-...} & sonnet-4-6 & 154 & readiness 90\%; 2,376,593 tokens \\
ablation & bloom\_v9c2\_ablation\_ndisc1       & \texttt{a551c74d-...} & sonnet-4-6 &  34 & n\_disc=1; all other v9c2 remediations ON; readiness\_failed (53\%); 631,243 tokens \\
\bottomrule
\end{tabular}
\end{table}

Full 36-character run UUIDs: v4 \texttt{19b39e20-6901-4c90-8f3c-e24b352f23ae};
v5 \texttt{02070b48-308f-4bf3-bac9-2cc76563ade9};
v6 \texttt{cd5d2f0e-8036-41d6-b755-b6c6ba08ab3e};
v7a \texttt{dd2fd888-0e7b-40a9-ad2a-abb6e9b8b6d1};
v7b \texttt{32ffaa49-2d26-4edd-9cf8-ecdf1d019c78};
v8 \texttt{ef107ef9-bee8-420e-a94f-31aee4a56c37};
v9b-smoke \texttt{9b599c12-ab43-42f8-8af4-de26551b3981};
v9b-prod \texttt{772d2ce7-246a-467e-89ae-b63d02d39c85};
v9c \texttt{b412cfdb-0522-4997-a47e-1738eb414f4b};
v9c2 \texttt{dc1bdd27-6ae0-46be-be74-fc8a974bb328};
ablation \texttt{a551c74d-aaed-498f-82cb-188b610380a8}.

\section{Confound $\to$ Remediation Map}
\label{appendix:confounds}

\begin{table}[h]
\centering
\caption{Confound taxonomy and remediations.}
\label{tab:confound-map}
\begin{tabular}{@{}llp{4cm}p{4.5cm}@{}}
\toprule
Confound & Description & Affected generations & Remediation \\
\midrule
F1  & Static corrective feedback           & v8--v9c2        & (open) \\
F2  & Exact-match scoring artifact (L6)    & v8, v9b, v9c    & A8: exclude L6 + semantic judge \\
F3  & Discussion context not injected      & v8 (partial), v9b, v9c & A0: inject learner memory \\
F4  & Pseudoreplication (unique\_discussions=1) & v4--v9c (v8 partial) & A3: population-multiplication, n\_disc=5 \\
F5a & Learner-profile homogeneity          & v4, v7a, v7b (scoped) & A4: 4 balanced types (satisfied v8+) \\
F5b & Single-instance role agents          & v4--v9c         & A5: per-discussion teacher instances \\
F5c & Learning-budget asymmetry (lecture-only) & v9b, v9c   & A6: self-reflection phase \\
\bottomrule
\end{tabular}
\end{table}

\section{Reproducibility}
\label{appendix:reproducibility}

Each generation has a version-pinned config
(\texttt{experiments/bloom\_1n\_vs\_1on1/config/config\_v*.yml}) and run script
(\texttt{scripts/run\_experiment\_v*.rb}).  Reports regenerate from a frozen run
database via \texttt{scripts/regen\_report.rb}.  Audit scripts
\texttt{check\_replication.py} (counts \texttt{unique\_discussions} per condition)
and \texttt{check\_confounds.py} (single-instance and profile-diversity checks)
reproduce the cross-version audit table.  The Zarn Token reference lecture,
evaluation tasks, and rubric are included under \texttt{domains/zarn\_tokens/}.

% -----------------------------------------------------------------------
\section*{Acknowledgement of Method}

This paper's negative result was only trustworthy because we audited our own
instrument before believing it.  We recommend the same discipline --- replication
counting, single-instance checks, prerequisite gating, and cost normalization --- as
a default for agent-based theory testing.

% -----------------------------------------------------------------------
\bibliographystyle{apalike}
\begin{thebibliography}{99}

\bibitem[Aher et al., 2023]{aher2023}
Aher, G., Arriaga, R.~I., \& Kalai, A.~T. (2023).
\newblock Using large language models to simulate multiple humans and replicate
  human subject studies.
\newblock \emph{Proceedings of the 40th International Conference on Machine
  Learning (ICML 2023)}, PMLR~202.

\bibitem[Argyle et al., 2023]{argyle2023}
Argyle, L.~P., Busby, E.~C., Fulda, N., Gubler, J.~R., Rytting, C., \&
  Wingate, D. (2023).
\newblock Out of one, many: Using language models to simulate human samples.
\newblock \emph{Political Analysis}, \textbf{31}(3), 337--351.
\newblock \url{https://doi.org/10.1017/pan.2023.2}

\bibitem[Bloom, 1984]{bloom1984}
Bloom, B.~S. (1984).
\newblock The 2 sigma problem: The search for methods of group instruction as
  effective as one-to-one tutoring.
\newblock \emph{Educational Researcher}, \textbf{13}(6), 4--16.
\newblock \url{https://doi.org/10.3102/0013189X013006004}

\bibitem[Hurlbert, 1984]{hurlbert1984}
Hurlbert, S.~H. (1984).
\newblock Pseudoreplication and the design of ecological field experiments.
\newblock \emph{Ecological Monographs}, \textbf{54}(2), 187--211.
\newblock \url{https://doi.org/10.2307/1942661}

\bibitem[Park et al., 2023]{park2023}
Park, J.~S., O'Brien, J.~C., Cai, C.~J., Morris, M.~R., Liang, P., \&
  Bernstein, M.~S. (2023).
\newblock Generative agents: Interactive simulacra of human behavior.
\newblock \emph{Proceedings of the 36th Annual ACM Symposium on User Interface
  Software and Technology (UIST 2023)}.
\newblock \url{https://doi.org/10.1145/3586183.3606763}

\bibitem[Rajpurkar et al., 2016]{rajpurkar2016}
Rajpurkar, P., Zhang, J., Lopyrev, K., \& Liang, P. (2016).
\newblock {SQuAD}: 100,000+ questions for machine comprehension of text.
\newblock \emph{Proceedings of the 2016 Conference on Empirical Methods in
  Natural Language Processing (EMNLP 2016)}, 2383--2392.
\newblock \url{https://doi.org/10.18653/v1/D16-2945}

\end{thebibliography}

\end{document}
